% Source: https://artofproblemsolving.com/community/c3h1707222_mathcounts_problems
% Stable identity: aops-topic-1707222-q03
Any five of the points form a convex pentagon (they lie on a circle), and the pentagon is determined by the two points left out. Label the points $0,1,\dots,6$ around the circle. Measure the sides of a pentagon by the number of steps between consecutive vertices; a side of $k$ steps is a chord of length $2\sin\dfrac{k\pi}{7}$, and congruent pentagons have the same cyclic sequence of sides up to rotation and reflection.

Rotating, we may leave out $0$ and $k$ with $k\in\{1,2,3\}$ (leaving out $0$ and $7-k$ is the mirror image).
\begin{itemize}
\item $k=1$ (omit $0,1$): sides $1,1,1,1,3$.
\item $k=2$ (omit $0,2$): vertices $1,3,4,5,6$, sides $2,1,1,1,2$: the two long sides are adjacent (meeting at vertex $1$).
\item $k=3$ (omit $0,3$): vertices $1,2,4,5,6$, sides $1,2,1,1,2$: the two long sides are not adjacent.
\end{itemize}
The first pentagon has a side of $3$ steps and the others do not; the last two have the same side lengths, but in a different cyclic order, so they are not congruent. Hence there are $\boxed{3}$ non-congruent pentagons.
